\chapter{Компоненты временных рядов2}
\label{chap: time_series_components2}

\section{Введение}

Временные ряды часто представляют собой сложную смесь различных влияний
, что затрудняет выявление истинных закономерностей и прогнозирование.
Чтобы за хаосом увидеть структуру, временной ряд принято разделять на
составляющие компоненты:
\begin{enumerate}
  \item Тренд (жирным) - показывает общее долгосрочное направление
    развития ряда.
  \item Сезонность (жирным) - отражает повторяющиеся колебания с
    фиксированной периодичностью.
  \item Ошибка- случайная непрогнозируемая компонента. Она может
    включать в себя несистематические шумы и аномалии.
  \item Цикл (жирным) - как и сезонность, показывает повторяющиеся
    подъемы и спады вокруг тренда.
    Однако ее ключевое отличие – отсутствие фиксированной периодичности.

    \ex{}{
    \item Краткосрочные (Китчина, 2-4 года): Связаны с колебаниями запасов.

    \item Среднесрочные (Жюгляра, 7-11 лет): Связаны с инвестициями в
      основной капитал.

    \item   Долгосрочные (Кузнеца, 15-25 лет): Связаны с демографией
      и миграцией.
    \item   Длинные волны (Кондратьева, 50-60 лет): Связаны с
      технологическими революциями.
    \item

    }

    Ключевая особенность (жирным): Циклы имеют нестрогую
    периодичность — их длительность и сила могут меняться под
    влиянием внешних факторов (кризисы, войны, технологии).

    Такой подход — декомпозиция — позволяет не просто глубже понять
    природу данных, но и значительно повысить точность прогнозов,
    планируя стратегии с учетом устойчивых тенденций и предсказуемых
    сезонных всплесков или спадов.
    По сути, это мощный метод превращения данных в ценную информацию
    для принятия взвешенных решений.

    \section{MECE декомпозиция}

    \label{sec: mece_decomposition}

    \ex{}{
      Если перед вами стоит задача предсказания прибыли компании,
      имеет смысл предварительно
      представить ее в виде разности выручки и затрат и далее
      отдельно прогнозировать
      каждую компоненту. Вспомним базовую микроэкономику:

      \[
        \begin{aligned}
          \text{Profit} &= \text{Total Revenue} - \text{Total Costs} \\
          \text{Total Revenue} &= \text{Average Check} \times
          \text{Quantity Sold} \\
          \text{Total Costs} &= \text{Variable Costs} + \text{Fixed Costs} \\
          \text{Variable Costs} &= \text{Average Variable Cost}
          \times \text{Quantity Sold} \\
          \text{Profit} &= (\text{Average Check} - \text{Average
          Variable Cost}) \times
          \text{Quantity Sold} - \text{Fixed Costs}
        \end{aligned}
      \]

      Разложив прибыль в итоговом виде, мы можем создать
      предсказательную модель для каждого элемента и затем получить
      итоговый прогноз прибыль, собрав в композицию прогнозы.
    }
